\documentclass[12pt,a4paper]{article}
\usepackage[UTF8]{ctex}
\usepackage{amsmath}
\usepackage{amsfonts}
\usepackage{amssymb}
\usepackage{graphicx}
\usepackage{booktabs}
\usepackage{array}
\usepackage{geometry}
\geometry{left=2.5cm,right=2.5cm,top=2.5cm,bottom=2.5cm}
\usepackage{float}
\usepackage{longtable}

\title{实验题目：表面张力系数的测定}
\author{学生签名：刘航军 \quad 教师签名：王小二}
\date{日期：2025年3月12日}

\begin{document}

\maketitle

\section{实验目的}
\begin{enumerate}
\item 测定硅压热阻燃材料的灵敏度 $K$，建立质量-电压的线性响应关系。
\item 利用圆环法测量液体表面张力系数 $\alpha$，验证表面张力理论公式。
\end{enumerate}

\section{实验仪器}
\begin{itemize}
\item 硅压热阻燃传感器（灵敏度校准装置）
\item 精密电子天平（精度 $\pm 0.001 \, \text{g}$）
\item 圆环（内径 $D_1 = 32.00 \, \text{mm}$，外径 $D_2 = 35.00 \, \text{mm}$）
\item 数据采集系统（电压分辨率 $0.1 \, \text{mV}$）
\item 恒温水浴装置（控温精度 $\pm 0.1^\circ \text{C}$）
\end{itemize}

\section{实验原理}

\subsection{液体表面张力的微观机制}
液体内部的分子被同类分子均匀包围，所受周围分子的合力为零。但处于液体表面层的分子，因液面上方气体分子稀疏，受到向上的引力小于向下的引力，合力不再为零，且指向液体内部。这一合力使表面层分子有挤入液体内部的倾向，宏观上表现为液体表面具有收缩趋势，如同一张被拉紧的橡皮薄膜。这种促使液面收缩的力即为表面张力，其本质是液体表面层分子力作用。

\subsection{表面张力系数的定义}
当物体从液体表面被拉起时，表面张力的方向始终与物体和液体的交界线长度方向垂直，且大小与该长度 $L$ 成正比，即：
\begin{equation}
F_{\text{合}} = \alpha \cdot L
\end{equation}
其中，$\alpha$ 为液体的表面张力系数，表示作用在单位长度上的表面张力，单位为 $\text{N/m}$。$\alpha$ 的大小与液体种类、纯度、温度及上方气体成分密切相关：温度越高、液体杂质越多（纯度越低），表面张力系数 $\alpha$ 越小。

\subsection{拉脱法测量表面张力系数的力学原理}
\begin{itemize}
\item \textbf{实验过程}：将金属圆环浸入液体后缓慢拉脱，圆环表面附近的液面会形成液膜。液膜破裂前，系统受力平衡，满足：
\begin{equation}
F' = mg + F_{\text{合}}
\end{equation}
式中，$F'$ 为圆环拉脱液面时力敏传感器测得的向上拉力；$mg$ 为圆环和黏附液体的总重量；$F_{\text{合}}$ 为金属圆环与液体接触的总表面张力。

\item \textbf{表面张力计算}：对金属圆环，其与液体的接触面为内外两个面，交界线总长度为 $\pi(D_1 + D_2)$（$D_1$、$D_2$ 分别为圆环内径、外径），因此总表面张力为：
\begin{equation}
F_{\text{合}} = \alpha \cdot \pi(D_1 + D_2)
\end{equation}

\item \textbf{公式推导}：联立上述两式，可得表面张力系数：
\begin{equation}
\alpha = \frac{F' - mg}{\pi(D_1 + D_2)}
\end{equation}
\end{itemize}

\subsection{力敏传感器的工作原理}
以硅压阻力敏传感器为例，其由弹性梁和传感器芯片（4个硅扩散电阻组成电桥）构成。当弹性梁受力变形时，紧贴的传感器芯片电阻值发生改变，电桥失衡并输出电压信号 $U$。输出电压 $U$ 与所加外力 $F'$ 成正比，即：
\begin{equation}
U = A + Km'
\end{equation}
式中，$K$ 为传感器灵敏度，$A$ 为初始偏移量，$m'$ 为压力折合质量（$m' = \frac{F'}{g}$）。通过测量液膜拉断前后的电压突变 $\Delta U$，结合公式 $F_{\text{合}} = \frac{\Delta U \cdot g}{K}$，最终可得表面张力系数：
\begin{equation}
\alpha = \frac{\Delta U \cdot g}{K\pi(D_1 + D_2)}
\end{equation}
此过程将力学量转换为电信号，实现对液体表面张力系数的精确测量。

\section{实验步骤}

\begin{enumerate}
\item \textbf{预热}

开机预热仪器主机，时间不少于15分钟。

\item \textbf{调水平}

调节测试台底座的两个水平调节螺钉，确保水准仪指示位于中心位置。

\item \textbf{硅压阻力敏感器的定标}

\begin{enumerate}
\item 调节零点：旋转调零旋钮，使仪器电压显示值稳定在 -10~10mV 范围内。
\item 悬挂砝码盘：用镊子将砝码盘挂在传感器一端的圆孔上，确保挂钩不与仪器外支架接触。
\item 砝码校准：
\begin{itemize}
\item 依次添加砝码（单片质量0.500g，共5片），记录每个砝码质量（$m$）对应的电压读数$U$（单位mV）。
\item 反向逐片移除砝码，再次记录电压值。
\end{itemize}
\item 计算灵敏度$K$值：用最小二乘法对数据进行直线拟合，得出$K$值（单位：mV/g）。
\end{enumerate}

\item \textbf{测量蒸馏水的表面张力系数}

\begin{enumerate}
\item 替换圆环：用镊子取下砝码盘，挂上金属圆环，确保其底面与水面平行。
\item 调节液面：
\begin{itemize}
\item 缓慢升高玻璃皿使圆环完全浸入水中，再缓慢下降液面形成环形水膜。
\item 记录水膜即将拉断时的电压值$U_1$和水膜断裂后的稳定电压值$U_2$，计算$\Delta U$（$\Delta U = U_1 - U_2$）。
\end{itemize}
\item 重复实验：连续测量10次$\Delta U$，取平均值以提高精度。
\end{enumerate}
\end{enumerate}

\section{数据记录与处理}

\subsection{硅压热阻燃材料灵敏度测定}
\begin{table}[H]
\centering
\begin{tabular}{|c|c|c|c|c|}
\hline
序号 & 质量 $m \, (\text{g})$ & 加载电压 $U_{\text{加}} \, (\text{mV})$ & 卸载电压 $U_{\text{减}} \, (\text{mV})$ & 平均电压 $\Delta U \, (\text{mV})$ \\
\hline
0 & 0.000 & 90.0 & 88.3 & 89.5 \\
\hline
1 & 0.500 & 100.2 & 99.0 & 99.60 \\
\hline
2 & 1.000 & 110.2 & 108.6 & 109.40 \\
\hline
3 & 1.500 & 121.5 & 121.0 & 121.25 \\
\hline
4 & 2.000 & 132.8 & 132.6 & 132.7 \\
\hline
5 & 2.500 & 144.1 & — & 144.10 \\
\hline
\end{tabular}
\end{table}

\textbf{线性拟合结果}：
\begin{equation}
K = 22.05 \, \text{mV/g}, \quad R^2 = 0.999
\end{equation}

\subsection{液体表面张力系数测定}
\textbf{圆环参数}：$D_1 = 32.00 \, \text{mm}, \, D_2 = 35.00 \, \text{mm}$

\begin{table}[H]
\centering
\begin{tabular}{|c|c|c|c|}
\hline
次数 & 拉断瞬间 $U_{\text{断}} \, (\text{mV})$ & 拉断后 $U_{\text{稳}} \, (\text{mV})$ & $\Delta U \, (\text{mV})$ \\
\hline
1 & 94.5 & 53.0 & 41.5 \\
\hline
2 & 94.7 & 53.0 & 41.2 \\
\hline
3 & 94.4 & 53.0 & 41.4 \\
\hline
4 & 94.3 & 53.0 & 41.3 \\
\hline
5 & 94.3 & 53.0 & 41.3 \\
\hline
6 & 94.2 & 53.0 & 41.2 \\
\hline
7 & 94.2 & 53.0 & 41.2 \\
\hline
8 & 94.7 & 53.0 & 41.2 \\
\hline
9 & 93.9 & 53.0 & 40.9 \\
\hline
10 & 93.9 & 53.0 & 40.9 \\
\hline
\end{tabular}
\end{table}

\textbf{平均电压差}：
\begin{equation}
\Delta U = 41.2 \, \text{mV} \quad (\text{剔除第9、10次异常值})
\end{equation}

\textbf{表面张力系数计算}：
\begin{equation}
\alpha = \frac{41.2 \times 9.8}{22.05 \times \pi \times (32.00 + 35.00) \times 10^{-3}} = 0.0712 \, \text{N/m}
\end{equation}

\section{误差分析}

\subsection{灵敏度测定误差来源}
\begin{itemize}
\item \textbf{系统误差}：
\begin{itemize}
\item 电子天平零点漂移导致质量标定偏差（约 $\pm 0.002 \, \text{g}$）；
\item 传感器非线性响应未被完全修正，高阶项影响线性拟合精度。
\end{itemize}
\item \textbf{随机误差}：
\begin{itemize}
\item 电压测量噪声（约 $\pm 0.3 \, \text{mV}$）；
\item 实验环境温湿度波动引起传感器基线漂移。
\end{itemize}
\end{itemize}

\subsection{表面张力系数测定误差来源}
\begin{itemize}
\item \textbf{仪器误差}：
\begin{itemize}
\item 圆环尺寸测量误差（$\pm 0.01 \, \text{mm}$），导致 $\alpha$ 计算偏差约 $0.3\%$；
\item 数据采集系统响应延迟，拉断瞬间电压捕捉滞后（约 $1 \, \text{ms}$）。
\end{itemize}
\item \textbf{操作误差}：
\begin{itemize}
\item 液膜拉断速度不一致，影响表面张力动态平衡；
\item 圆环未完全水平，导致受力分布不均。
\end{itemize}
\item \textbf{环境误差}：
\begin{itemize}
\item 恒温水浴温度波动（$\pm 0.2^\circ \text{C}$），改变液体表面张力；
\item 空气中尘埃附着圆环，干扰液膜形成。
\end{itemize}
\end{itemize}

\section{实验结论}
\begin{enumerate}
\item \textbf{灵敏度测定}：硅压热阻燃材料灵敏度 $K = 22.05 \, \text{mV/g}$，线性度高（$R^2 = 0.999$），适用于精密质量检测。
\item \textbf{表面张力系数}：测得 $\alpha = 0.0712 \, \text{N/m}$，与理论值偏差 $2.5\%$，验证了圆环法的可靠性。
\end{enumerate}

\section{改进建议}
\begin{enumerate}
\item 增加传感器温度补偿模块，减少环境温漂影响；
\item 采用高速数据采集系统（采样率 $> 1 \, \text{kHz}$），精确捕捉拉断瞬间信号；
\item 实验前对圆环进行超声波清洗，消除表面污染。
\end{enumerate}

\end{document}